\section{Related work}\label{sec:related}

\paragraph{Object detection.} One-stage detectors such as YOLO~\cite{redmon2016yolo} and its Ultralytics successor YOLOv8~\cite{jocher2023yolov8}, trained on COCO~\cite{lin2014coco}, predict boxes and classes in a single network pass. They process frames independently and do not exploit temporal redundancy. We use YOLOv8m as a black box; nothing in the method depends on the particular detector.

\paragraph{Exploiting temporal redundancy.} One line of work reuses computation inside the network: Deep Feature Flow~\cite{zhu2017dff} computes expensive features on sparse key frames and propagates them with optical flow. Another line filters frames in front of the network. NoScope~\cite{kang2017noscope} cascades difference detectors and specialised models for binary video queries, Chameleon~\cite{jiang2018chameleon} adapts frame rate and resolution over time, Reducto~\cite{li2020reducto} filters frames on the camera using low-level difference features with query-specific thresholds, and FrameHopper~\cite{arefeen2022framehopper} learns a frame-skipping policy by reinforcement learning. These systems target multi-stream or query-driven analytics and optimise end-to-end query accuracy. We isolate the single-stream case, make the trade-off between cost and staleness explicit with provable bounds, and characterise the physical failure modes of the change detector itself. Tracking-by-detection, for example SORT~\cite{bewley2016sort}, is complementary: a tracker could extrapolate boxes between detector calls instead of retaining them.

\paragraph{Change detection.} Frame differencing is the oldest change detector~\cite{jain1979difference}. Adaptive Gaussian mixtures~\cite{stauffer1999adaptive}, available in OpenCV~\cite{bradski2000opencv} as MOG2~\cite{zivkovic2004improved}, maintain a per-pixel background model. The survey of Garcia-Garcia et al.~\cite{garcia2020background} lists illumination changes, camera jitter and noise among the main open challenges of background subtraction, which are precisely the factors analysed here. We use MOG2 as a stronger but costlier baseline gate.

\paragraph{Image registration.} Global translations are classically estimated by phase correlation~\cite{kuglin1975phase}, which is non-iterative and insensitive to global intensity changes. Iterative methods such as Lucas--Kanade~\cite{lucas1981iterative} or enhanced-correlation-coefficient maximisation~\cite{evangelidis2008ecc} are more accurate for sub-pixel motion but more expensive. The stabilised gate uses phase correlation and guards it with a model-selection test.
